% Obsahový zdroj: import_hodiny/41_Smisene_kombinatoricke_ulohy.tex a odpovídající PDF (obě strany).
% P1–P6: zdrojové Příklady 1–6; C1–C6: zdrojové Procvičování 1–6.
% Matematický obsah zachován; výklad, odvození a řešení pedagogicky rozvedeny.
% Upřesnění a opravy: viz KMB_29-41_Kontrola.md; původní import zůstává beze změny.
% Rozpis: 4_rocnik/tematicke_plany/Rozpis_vyuky_MAT_4_rocnik.tex, hodina 41.
% Plán: 4_rocnik/tematicke_plany/SPS.pdf, oddíly 2.1–2.3.
% Sazba: sablony/pripravy/TEMPLATE_Priprava.tex a společný priprava.sty.
% Vlastní doplnění: didaktický výklad, zdůvodnění, metodika; žádné nové číslované úlohy.
% Učebnice/09_Kombinatorika nemá obsahovou kapitolu; místní kombinatorické listy ani Realisticky.cz nejsou dostupné.
\documentclass[11pt,a4paper]{article}
% Spolecny styl pro ucitelske pripravy MAT (Work/MAT)
% Verze 1.1, 2026-10-23
\NeedsTeXFormat{LaTeX2e}
\ProvidesPackage{priprava}[2026/10/23 v1.1 Sablona ucitelskych priprav MAT]

% Zakladni typografie a cestina
\RequirePackage[T1]{fontenc}
\RequirePackage[utf8]{inputenc}
\RequirePackage[czech]{babel}
\RequirePackage{lmodern}
\RequirePackage{microtype}
\RequirePackage{amsmath,amssymb}
\RequirePackage{xcolor}
\RequirePackage{enumitem}
\RequirePackage[most]{tcolorbox}
\RequirePackage{tabularx,booktabs,array}
\RequirePackage{fancyhdr}
\RequirePackage{needspace}
\RequirePackage{geometry}
\RequirePackage[hidelinks]{hyperref}

\geometry{a4paper,top=17mm,bottom=18mm,left=23mm,right=23mm}
\definecolor{notegray}{gray}{0.39}
\setlength{\parindent}{0pt}
\setlength{\parskip}{2.5pt plus 1pt minus 0.5pt}
\setlist[itemize]{topsep=2pt,itemsep=1pt,parsep=0pt,leftmargin=1.45em}
\setlist[enumerate]{topsep=2pt,itemsep=2pt,parsep=0pt,leftmargin=1.65em}
\setcounter{secnumdepth}{0}
\renewcommand{\arraystretch}{1.12}

% Zahlavi: vlevo rocnik | globalni tema, vpravo cislo: nazev hodiny
% Zapati zustava zcela prazdne.
\pagestyle{fancy}
\fancyhf{}
\renewcommand{\headrulewidth}{0.35pt}
\renewcommand{\footrulewidth}{0pt}
\setlength{\headheight}{14pt}
\newcommand{\nastavPripravu}[4]{%
  \fancyhead[L]{\small #1\enspace\textbar\enspace #2}%
  \fancyhead[R]{\small #3: #4}%
}

% Minimalisticke cernobile boxy. Kazda definice zvlast.
\tcbset{commonlesson/.style={enhanced,breakable=false,colback=white,colframe=black,
  colbacktitle=white,coltitle=black,boxrule=0.45pt,
  arc=0.8mm,outer arc=0.8mm,left=3mm,right=3mm,top=1.6mm,bottom=1.6mm,
  toptitle=1mm,bottomtitle=0.7mm,fonttitle=\bfseries,
  attach boxed title to top left={xshift=2.5mm,yshift=-1.15mm},
  boxed title style={colback=white,colframe=white,boxrule=0pt,left=1mm,right=1mm}}}
\newtcolorbox{zadaniBox}[1][]{commonlesson,title={Zadání},#1}
\newtcolorbox{prikladBox}[1][]{commonlesson,title={Příklad},#1}
\newtcolorbox{definiceBox}[1][]{commonlesson,title={Definice},#1}

% Spolecny box pro zadani a jeho reseni, bez mezery mezi castmi.
% Pouziti:
%   \begin{prikladReseniBox}
%     Zadani ulohy.
%     \tcblower
%     \begin{align*} ... \end{align*}
%   \end{prikladReseniBox}
% Spodni cast ma jemny sedy podklad a delici linku; lze v ni sazet
% vice radku matematiky. Obe casti maji jen jeden spolecny ramecek.
\tcbset{solutionlesson/.style={
  commonlesson,
  enhanced jigsaw,
  breakable,
  colbacklower=black!3,
  segmentation style={solid,draw=black!25,line width=0.35pt},
  fontlower=\small\color{notegray},
  before lower={\noindent\textbf{Řešení:}\enspace},
  leftlower=3mm,rightlower=3mm
}}
\newtcolorbox{prikladReseniBox}[1][]{solutionlesson,title={Příklad},#1}
\newtcolorbox{zadaniReseniBox}[1][]{solutionlesson,title={Zadání},#1}

% Pomucky pro autora pripravy: sedive poznamky; \reseni pro starsi pripravy
\newcommand{\poznamka}[1]{\par\begingroup\small\color{notegray}\itshape
  \textnormal{\textbf{Poznámka pro učitele.}}\enspace #1\par\endgroup}
\newcommand{\reseni}[1]{\par\begingroup\small\color{notegray}
  \textbf{Očekávané řešení:}\enspace #1\par\endgroup}
\newcommand{\blok}[1]{\needspace{5\baselineskip}\subsection*{#1}}
\newcommand{\vyrok}[1]{\uv{#1}}

\nastavPripravu{4. ročník}{Kombinatorika}{41}{Smíšené kombinatorické úlohy}
\begin{document}


Po průběžném testu vrať do středu pozornosti volbu modelu. Nech žáky u každé úlohy ještě před výpočtem vysvětlit, co tvoří jeden výsledek a proč nestačí bez rozmyslu dosadit do jednoho vzorce.
\poznamka{Předpoklady: hodiny 29–39; zejména součet a součin, variace, permutace, kombinace a doplňková metoda. Cíl: žák zdůvodní model slovní úlohy s omezením, rozčlení úplné disjunktní případy a zkontroluje výsledek. Jádro tvoří P1 a P3; ostatní příklady jsou zásoba pro rozbor chyb z testu a diferenciaci. Není účelem vypočítat všechna zadání v jedné hodině.}

\blok{Výklad a odvození}

\textbf{Postup řešení.} Nejdříve popiš jeden výsledek, rozliš pořadí a opakování a urči počet dostupných i vybíraných prvků. Potom přečti všechna omezení. Teprve nyní zvol vzorec nebo sestav součin dílčích voleb. Je-li třeba více případů, musí se vzájemně vylučovat a společně pokrýt všechny přípustné výsledky.

\begin{tabularx}{\linewidth}{@{}>{\raggedright\arraybackslash}p{45mm}X@{}}
\toprule
\textbf{Situace} & \textbf{Nástroj a podmínky}\\\midrule
Řazení všech různých prvků & $P(n)=n!$, celé $n\geq0$.\\
Obsazení $k$ různých rolí & $V(n,k)=n!/(n-k)!$, celá $0\leq k\leq n$.\\
Volné opakování v kódu & $\overline V(n,k)=n^k$, celá $n\geq1$, $k\geq0$.\\
Pevné četnosti stejných kopií & $P(n;n_1,\ldots,n_r)=n!/(n_1!\cdots n_r!)$, nezáporné celé četnosti se součtem $n$.\\
Skupina bez rozlišených rolí & $C(n,k)=\binom nk$, celá $0\leq k\leq n$, bez opakování.\\
\bottomrule
\end{tabularx}

Slova „právě jeden“ a „alespoň jeden“ neznamenají totéž. U „alespoň jednoho“ často pomůže doplněk „žádný“. Požadavek sousedství vede k bloku a ke kontrole jeho vnitřního pořadí. Číselné úlohy s nulou vyžadují odlišit kód od čísla a někdy rozdělit případy podle poslední cifry.

Pravidlo součinu nevyžaduje stejnou nabídku znaků v každé větvi: vyžaduje stejný \emph{počet} pokračování příslušného kroku po každé předchozí přípustné volbě. Například volba poslední nuly a poslední nenulové sudé cifry může měnit počet možností na začátku; pak se součiny těchto disjunktních případů sčítají.


\poznamka{Průběh: 0–7 min rozhodovací otázky a návrat k chybám z testu; 7–24 min P1 a P3; 24–36 min samostatně C1–C2 a jedna z C4–C6 podle zjištěných potřeb; 36–43 min porovnání metod, případně krátký návrat B1; 43–45 min návaznost na výsledky náhodných pokusů. P2, P4–P6 a zbývající C úlohy jsou rezerva či domácí práce. Obsah je záměrně širší než 45 minut, aby učitel mohl volit podle testu.}

\blok{Společně řešené příklady}

% Zdroj P1: import_hodiny/41_Smisene_kombinatoricke_ulohy.tex, Příklad 1.
\begin{prikladReseniBox}[title={Příklad P1}]
Kolik tříčlenných komisí bez rozlišených funkcí lze vybrat z devíti osob (pěti žen a čtyř mužů), má-li být v komisi alespoň jedna žena?
\tcblower
Všech komisí je $\binom93=84$. Nevyhovují právě komise bez ženy, tedy tři muži ze čtyř: $\binom43=4$. Doplňkem dostaneme $84-4=80$. Kontrola disjunktními případy: právě jedna žena $5\binom42=30$, dvě ženy $\binom52\cdot4=40$, tři ženy $\binom53=10$. Součet $30+40+10=80$.
\end{prikladReseniBox}


% Zdroj P2: import_hodiny/41_Smisene_kombinatoricke_ulohy.tex, Příklad 2.
\begin{prikladReseniBox}[title={Příklad P2}]
Kolik pětiznakových kódů ze čtyř různých symbolů, mezi nimiž je A, obsahuje alespoň jedno A? Opakování je dovoleno.
\tcblower
Celkem \(4^5\) kódů; bez A je \(3^5\). Počet \(4^5-3^5=1024-243=781\).
\end{prikladReseniBox}


% Zdroj P3: import_hodiny/41_Smisene_kombinatoricke_ulohy.tex, Příklad 3.
\begin{prikladReseniBox}[title={Příklad P3}]
Kolik čtyřciferných sudých čísel bez opakování vytvoříme z číslic $0,1,2,3,4,5,6$?
\tcblower
Sudé číslo končí $0,2,4$ nebo $6$. Tyto případy rozdělíme podle vlivu nuly na první pozici.
\begin{itemize}
\item Konec $0$: první cifra má šest nenulových možností, prostřední dvě pak pět a čtyři. Dostáváme $6\cdot5\cdot4=120$.
\item Konec $2,4$ nebo $6$: tři možnosti konce; první cifra má pět zbývajících nenulových možností. Po obsazení obou krajů je pro druhou pozici pět zbývajících cifer včetně nuly a pro třetí čtyři. Dostáváme $3\cdot5\cdot5\cdot4=300$.
\end{itemize}
Případy jsou disjunktní a vyčerpávají všechna sudá zakončení. Celkem $120+300=420$ čísel. Jeden stejný součin pro všechny čtyři konce by nezohlednil rozdílnou nabídku první cifry.
\end{prikladReseniBox}


% Zdroj P4: import_hodiny/41_Smisene_kombinatoricke_ulohy.tex, Příklad 4.
\begin{prikladReseniBox}[title={Příklad P4}]
Šest různých žáků v řadě, A a B nesmějí stát vedle sebe.
\tcblower
Všechna pořadí \(6!=720\). Nepřípustná: blok AB/BA plus čtyři lidé, \(2\cdot5!=240\). Přípustná: \(720-240=480\).
\end{prikladReseniBox}


% Zdroj P5: import_hodiny/41_Smisene_kombinatoricke_ulohy.tex, Příklad 5.
\begin{prikladReseniBox}[title={Příklad P5}]
Kolik osmimístných binárních řetězců obsahuje právě tři jedničky?
\tcblower
Vybereme tři pozice pro jedničky, ostatní jsou nuly. \(\binom83=56\). Stejný výsledek dává \(\frac{8!}{3!5!}\).
\end{prikladReseniBox}


% Zdroj P6: import_hodiny/41_Smisene_kombinatoricke_ulohy.tex, Příklad 6.
\begin{prikladReseniBox}[title={Příklad P6}]
Kolik tříčlenných skupin bez rozlišených funkcí lze vybrat ze sedmi lidí, nesmějí-li být předem určení lidé P a Q vybráni současně?
\tcblower
Všech trojic je $\binom73=35$. Zakázaná skupina obsahuje P i Q a jednoho z ostatních pěti, tedy $\binom51=5$ možností. Doplňkem $35-5=30$. Kontrola: ani P, ani Q dává $\binom53=10$ skupin; právě jeden z nich dává $2\binom52=20$. Součet je opět $30$.
\end{prikladReseniBox}


\newpage
\blok{Samostatná práce a zásoba k procvičení}

U každé úlohy nejprve napiš jednu větu zdůvodňující model. Při kontrole porovnej dvě různé metody, jsou-li dostupné. C3 předpokládá binární abecedu, C6 jednotkové kroky doprava a nahoru.

% Zdroj C1: import_hodiny/41_Smisene_kombinatoricke_ulohy.tex, Procvičování, položka 1.
\begin{zadaniReseniBox}[title={Úloha C1}]
Kolik tříčlenných skupin lze vybrat z osmi různých lidí, má-li být vybrán alespoň jeden ze dvou předem určených A, B? Funkce nerozlišujeme.
\tcblower
Všech skupin je $\binom83=56$. Bez A i B vybíráme ze šesti ostatních: $\binom63=20$. Vyhovuje $56-20=36$. Kontrola: právě jeden z A, B dává $2\binom62=30$, oba dávají šest voleb třetího člena. Celkem $30+6=36$.
\end{zadaniReseniBox}

% Zdroj C2: import_hodiny/41_Smisene_kombinatoricke_ulohy.tex, Procvičování, položka 2.
\begin{zadaniReseniBox}[title={Úloha C2}]
Kolik trojznakových kódů z pěti různých symbolů, mezi nimiž je X, obsahuje alespoň jedno X? Opakování je dovoleno.
\tcblower
Všech kódů je $5^3=125$, bez X jsou $4^3=64$. Doplňkem dostaneme $125-64=61$. Případy „X na první/druhé/třetí pozici“ nelze bez korekce sčítat, protože se překrývají.
\end{zadaniReseniBox}

% Zdroj C3: import_hodiny/41_Smisene_kombinatoricke_ulohy.tex, Procvičování, položka 3.
\begin{zadaniReseniBox}[title={Úloha C3}]
Kolik pětiznakových řetězců z abecedy $\{0,1\}$ obsahuje přesně dvě jedničky?
\tcblower
Vybereme dvě pozice z pěti pro jedničky. Zbylá tři místa jsou jednoznačně nuly. Počet $\binom52=10$, shodně $5!/(2!\,3!)=10$.
\end{zadaniReseniBox}

% Zdroj C4: import_hodiny/41_Smisene_kombinatoricke_ulohy.tex, Procvičování, položka 4.
\begin{zadaniReseniBox}[title={Úloha C4}]
Kolika způsoby lze na polici seřadit pět různých knih, musejí-li dvě předem určené knihy sousedit?
\tcblower
Dvě knihy spojíme do bloku. S dalšími třemi knihami řadíme čtyři různé objekty: $4!$ možností. Uvnitř bloku jsou dvě pořadí, proto $4!\cdot2=48$.
\end{zadaniReseniBox}

% Zdroj C5: import_hodiny/41_Smisene_kombinatoricke_ulohy.tex, Procvičování, položka 5.
\begin{zadaniReseniBox}[title={Úloha C5}]
Kolik čtyřčlenných skupin lze vybrat ze šesti žen a čtyř mužů, mají-li být právě dvě vybrané osoby ženy?
\tcblower
Vybereme dvě ženy $\binom62=15$ způsoby a dva muže $\binom42=6$ způsoby. Každá dvojice dílčích výběrů tvoří právě jednu skupinu, tedy $15\cdot6=90$.
\end{zadaniReseniBox}

% Zdroj C6: import_hodiny/41_Smisene_kombinatoricke_ulohy.tex, Procvičování, položka 6.
\begin{zadaniReseniBox}[title={Úloha C6}]
Kolik nejkratších cest po jednotkové mřížce vede z $(0,0)$ do $(3,3)$ pouze kroky doprava a nahoru?
\tcblower
Potřebujeme tři kroky doprava a tři nahoru. Vybereme tři ze šesti pozic pro kroky doprava: $\binom63=20$. Stejně jde o permutace kroků s opakováním $6!/(3!\,3!)=20$.
\end{zadaniReseniBox}


\blok{Krátký návrat k binomické větě (volitelně)}
% Zdroj B1: import_hodiny/39_Binomicka_veta.tex, Procvičování, položka 5.
% Totožná úloha je v této sadě v hodině 39 jako C5; nejde o nové zadání.
\begin{zadaniReseniBox}[title={Úloha B1}]
Najdi koeficient u $x^3$ v $(1-2x)^5$.
\tcblower
Obecný člen je $\binom5k1^{5-k}(-2x)^k$. Aby obsahoval $x^3$, musí být $k=3$; koeficient je $\binom53(-2)^3=-80$. Pořadí pěti závorek je pevné, vybíráme tři pozice, v nichž se uplatní $-2x$. Tím se binomická věta vrací k výběru pozic.
\end{zadaniReseniBox}

\poznamka{Typické chyby: započítání výběrů s A i B do dvou překrývajících se případů; ztráta zákazu počáteční nuly; násobení pořadí osob v nerozlišené komisi; sousedství řešené bez obou orientací. U P3 zkontroluj, že $120+300$ zahrnuje všechna dovolená sudá zakončení právě jednou. U B1 záporné znaménko patří do třetí mocniny.}

\textbf{Návaznost na hodinu 42.} Kombinatorika určila počet všech výsledků a počet výsledků splňujících podmínku. U náhodného pokusu bude navíc třeba přesně určit, co je jeden elementární výsledek a zda jsou všechny tyto výsledky stejně pravděpodobné. Samotné určení počtů ještě rovnoměrnost neprokazuje. V této hodině nezavádíme pravděpodobnostní vzorec; v další začneme výpisem výsledků hodů mincí a kostkou.
\end{document}
